\section*{历史注记（第五章与第六章）}
\phantomsection
\addcontentsline{toc}{section}{历史注记（第五章与第六章）}

（注意：罗马数字指本注记末尾所附的参考文献。）

不同阶的“无穷小”（或“无穷大”）之间的区分隐含地出现在微分学最早的著作中，例如费马的著作中；随着牛顿与莱布尼茨以及“高阶差分”理论的出现，它变得完全明确；而且人们很快就观察到，在最简单的情形，形如 \(f(x)/g(x)\) 的表达式在 f 与 g 都趋于 0 的点处的极限（或“真值”）由这些函数在所考虑点的邻域上的泰勒展开给出（“洛必达法则”，很可能属于约翰·伯努利）。

除初等情形外，从 XVII \(^{th}\) 世纪末起摆在数学家面前的主要的“渐近估计”问题，是当 n 很大时形如 \(\sum_{k=1}^{n} f(k)\) 的和的精确或近似计算；这种计算无论对插值法、对级数和的数值估计，还是对概率计算（其中诸如 n! 或 \(\binom{a}{n}\) 这类“大数的函数”起着占优的作用）都确实必不可少。早在牛顿，为了在 n 大时得到 \(\sum_{k=1}^{n} \frac{1}{a+k}\) 的近似值，就给出了一个方法，它（在这种特殊情形）归结为欧拉–麦克劳林公式 (I) 中前几项的计算。在该世纪末，雅各布·伯努利在概率计算的研究过程中着手确定和 \(S_{k}(n) = \sum_{p=1}^{n-1} p^{k}\) ——n 的一些多项式——并发现了其构成的一般规律（未给出证明）\(^{4}\) ，从而第一次在这些多项式系数的表达式中引入了以他的名字命名的数，以及用来计算它们的递推关系（(II), p.~97）。1730 年，斯特林对 \(\sum_{k=1}^{n} \log(x + ka)\) 得到了当 n 无限增大时的一个渐近展开，其办法是递归地计算系数。

欧拉关于级数及相关问题的决定性工作属于 1730 至 1745 年。置 \(\mathrm{S}(n) = \sum_{k=1}^{n} f(k)\) ，他把泰勒公式应用于函数 \(\mathrm{S}(n)\) ，得到

\[
f (n) = \mathrm{S} (n) - \mathrm{S} (n - 1) = \frac {d \mathrm{S}}{d n} - \frac {1}{2 !} \frac {d ^ {2} \mathrm{S}}{d n ^ {2}} + \frac {1}{3 !} \frac {d ^ {3} \mathrm{S}}{d n ^ {3}} - \dots ,
\]

对所得的方程，他用待定系数法加以“反演”，寻求形如下式的解

\[
\mathrm{S} (n) = \alpha \int f (n) d n + \beta f (n) + \gamma \frac {d f}{d n} + \delta \frac {d ^ {2} f}{d n ^ {2}} + \dots ;
\]

于是他逐项得到

\[
\mathrm{S} (n) = \int f (n) d n + \frac {f (n)}{2} + \frac {1}{1 2} \frac {d f}{d n} - \frac {1}{7 2 0} \frac {d ^ {3} f}{d n ^ {3}} + \frac {1}{3 0 2 4 0} \frac {d ^ {5} f}{d n ^ {5}} - \dots
\]

起初他并不能确定系数的构成规律（III a 与 d）。但在 1735 年前后，与多项式分解为线性因子的做法类比，他毫不犹豫地写下了公式

\[
\begin{array}{c} 1 - \frac {\sin s}{\sin \alpha} = \left(1 - \frac {s}{\alpha}\right) \left(1 - \frac {s}{\pi - \alpha}\right) \\ \left(1 - \frac {s}{- \pi - \alpha}\right) \left(1 - \frac {s}{2 \pi - \alpha}\right) \left(1 - \frac {s}{- 2 \pi - \alpha}\right) \dots \end{array}
\]

把两边按幂级数展开并令系数相等，他特别地（对 \(\alpha = \pi/2\) ）得到了级数 \(\sum_{k=1}^{\infty} \frac{1}{n^{2k}}\) 用 \(\pi\) 的幂表示的著名表达式（III b）\(^{5}\) 。几年之后，他终于看出 \(\pi\) 的这些幂的系数由与他的求和公式相同的方程给出，并认识到它们与伯努利引入的数、与 \(z/(e^{z} - 1)\) 的级数展开的系数之间的联系（III g）。

与欧拉相独立，麦克劳林在大约同一时间以稍欠冒险、接近于我们在正文中遵循的途径得到了同一个求和公式；他实际上迭代了“泰勒”公式，该公式把 \(f(x)\) 用差 \(f^{(2k+1)}(x+1)-f^{(2k+1)}(x)\) 表出，而这个公式是他用待定系数法“反演”这些差的泰勒展开而得到的（IV）；但他没有看出由欧拉发现的系数构成规律。

但麦克劳林与欧拉以及他那个时代的所有数学家一样，把所有这些公式都作为级数展开来陈述，其收敛性甚至未加研究。这并不是说收敛级数的概念在这个时期完全被忽视；自雅各布·伯努利以来人们就知道调和级数发散，欧拉甚至还借助他的求和公式估计了这个级数前 n 项的和，使这一结果更加明确（III c 与 d）；也正是欧拉指出，两个相邻伯努利数之比无限增大，因而以这些数为系数的幂级数不可能收敛（(III f), p.~357）\(^{6}\) 。但形式计算的倾向更为强大，欧拉超凡的直觉甚至也没能阻止他陷入荒谬，

例如他写出 \(0 = \sum_{n=-\infty}^{+\infty} x^n ((\text{III } f), \text{ p. 362})^7\) 时便是如此。

我们已经看到（第四章的历史注记），XIX \(^{th}\) 世纪初的数学家如何在对这种无节制、无根据的形式主义的厌倦中，把分析拉回到严格的道路上。一旦收敛级数的概念得到精确化，就需要一些简单的判别法，以便通过与已知的积分或级数比较来证明积分和级数的收敛；柯西在他的《代数分析》（Va）中给出了若干这样的判别法，而阿贝尔在一篇遗著（VI）中得到了对数收敛判别法。另一方面，柯西（V b）阐明了诸如斯特林级数这类（由应用欧拉–麦克劳林公式得到的、常称为“半收敛级数”的）级数的悖论；他证明，如果（鉴于欧拉关于伯努利数的评注）这类级数的通项 \(u_{k}(n)\) 当 n 取固定值时随 k 无限增大，那么尽管如此，当 k 取固定值时，部分和 \(s_{k}(n) = \sum_{h=1}^{k} u_{h}(n)\) 仍给出级数所“表示”的函数的一个渐近展开（当 n 趋于 \(+\infty\) 时），且 k 越大越精确。

在经典分析的大多数计算中，能够得到函数渐近展开的项数任意多的一般的构成规律；这一事实助长了级数与渐近展开之间（至少在用语上）经久不衰的混淆；以致 H. 庞加莱在 1886 年（VIII）着手把渐近展开的初等规则（以 \(1 / x\) 的整数幂为尺度、在 \(+\infty\) 的邻域上）整理成典时，仍然使用级数论的词汇。只是随着来自解析数论的渐近展开的出现，渐近展开与级数这两个概念之间的明确区分才最终建立，原因在于：在这门理论处理的大多数问题中，人们只能显式地得到所求展开的非常少的几项（多数情形只有一项）。

这些问题还使数学家习惯于使用变量的（实或整数）幂以外的比较尺度。这一推广首先应归功于杜·布瓦–雷蒙的工作（VII），他第一个系统地处理了函数在一点的邻域中的比较，并在极富独创性的工作中认识到比较尺度的“非阿基米德”特征，同时一般地研究了比较关系的积分与求导，并导出大量有趣的推论（VII b）。他的证明有时缺乏清晰性与严格性，而杜·布瓦–雷蒙结果的一个正确陈述应归功于哈代（IX）：他的主要贡献在于认识并证明了一组“初等函数”——(H) 函数——的存在，在这组函数上，分析的通常运算（特别是求导）可以应用于比较关系 \({}^{8}\) 。

